% ================= Topic 3 (VI) -- part b: problems 16--30 =================
\providecommand{\cpTooth}[3]{% #1 rows, #2 columns, #3 unit (cm): grid of toothpicks
  \foreach \j in {0,...,#1}{\foreach \i in {1,...,#2}{
    \draw[tmOrange!85!black,line width=1.6pt,line cap=round]
      ({(\i-1)*#3+0.05},{\j*#3})--({\i*#3-0.05},{\j*#3});}}
  \foreach \j in {1,...,#1}{\foreach \i in {0,...,#2}{
    \draw[tmOrange!85!black,line width=1.6pt,line cap=round]
      ({\i*#3},{(\j-1)*#3+0.05})--({\i*#3},{\j*#3-0.05});}}}

\begin{probbox}{16}{\relax}
Trên lưới điểm ở hình bên, bạn Jane đứng ở điểm $A$. Jane tung một đồng xu cân đối để quyết định
cách di chuyển: nếu đồng xu ra mặt \emph{sấp} thì Jane đi \emph{lên} một điểm; nếu ra mặt \emph{ngửa}
thì Jane đi \emph{sang phải} một điểm. Sau đúng bốn lần tung, Jane sẽ dừng ở một trong năm điểm
$P$, $Q$, $R$, $S$, $T$. Xác suất để Jane dừng ở điểm $R$ là bao nhiêu?
\begin{figC}\begin{scope}[scale=0.8]
  \foreach \y in {0,...,4}{\pgfmathtruncatemacro\m{4-\y}
    \foreach \x in {0,...,\m}{\fill[Navy] (\x,\y) circle (2.6pt);}}
  \node[below left] at (0,0) {$A$};
  \node[above right] at (0,4) {$P$};\node[above right] at (1,3) {$Q$};
  \node[above right] at (2,2) {$R$};\node[above right] at (3,1) {$S$};
  \node[above right] at (4,0) {$T$};
  \fill[tmOrange] (0,0) circle (3.4pt);
\end{scope}\end{figC}
\choices{$\dfrac12$}{$\dfrac38$}{$\dfrac9{16}$}{$\dfrac7{16}$}{$\dfrac5{16}$}
\end{probbox}
\probsol{
\textbf{Phân tích đề.} Mỗi lần tung, Jane đi \emph{lên} ($\uparrow$, khi ra sấp) hoặc \emph{sang phải}
($\rightarrow$, khi ra ngửa). Sau $4$ lần tung, Jane đi đúng $4$ bước. Ta cần tính xác suất để
sau $4$ bước Jane đứng ở $R$.

\step{1}{Đếm tất cả các kết quả có thể}
Mỗi lần tung có $2$ kết quả (sấp hoặc ngửa), bốn lần tung liên tiếp cho
$2\times2\times2\times2=16$ dãy kết quả, ví dụ $\uparrow\rightarrow\rightarrow\uparrow$.
Vì đồng xu cân đối nên $16$ dãy này \emph{đồng khả năng}.
\begin{recallbox}\emph{Quy tắc nhân:} nếu việc thứ nhất có $a$ cách, việc thứ hai có $b$ cách, \dots{} thì
làm lần lượt các việc đó có $a\times b\times\cdots$ cách. \emph{Xác suất} của một biến cố (khi các kết
quả đồng khả năng) $=\dfrac{\text{số kết quả thuận lợi}}{\text{số kết quả có thể}}$.\end{recallbox}

\step{2}{Điểm cuối chỉ phụ thuộc vào số bước đi lên}
Nếu trong $4$ bước có $k$ bước đi lên thì có $4-k$ bước sang phải. Thứ tự các bước không làm thay
đổi điểm cuối. Đặt $A$ là gốc, ta có: $4$ lần lên $\to P$; $3$ lần lên $\to Q$;
$2$ lần lên $\to R$; $1$ lần lên $\to S$; $0$ lần lên $\to T$.
Vậy Jane dừng ở $R$ khi và chỉ khi có \textbf{đúng $2$ bước lên và $2$ bước sang phải}.

\step{3}{Liệt kê các dãy thuận lợi}
Ta chọn vị trí của hai bước $\uparrow$ trong bốn lượt (lượt 1, 2, 3, 4):
\[\uparrow\uparrow\rightarrow\rightarrow,\quad \uparrow\rightarrow\uparrow\rightarrow,\quad
\uparrow\rightarrow\rightarrow\uparrow,\quad \rightarrow\uparrow\uparrow\rightarrow,\quad
\rightarrow\uparrow\rightarrow\uparrow,\quad \rightarrow\rightarrow\uparrow\uparrow.\]
Có $6$ dãy (chọn $2$ lượt trong $4$ lượt: $\dfrac{4\times3}{2}=6$ cách).

\step{4}{Tính xác suất}
\[P(\text{dừng ở } R)=\frac{6}{16}=\frac38.\]

\step{5}{Một cách nhìn bằng hình: ghi số đường đi tới mỗi điểm}
Ghi vào mỗi điểm số dãy bước dẫn từ $A$ tới điểm đó. Mỗi điểm chỉ có thể tới từ điểm ngay bên trái hoặc
ngay bên dưới, nên số ở mỗi điểm $=$ tổng hai số ở bên trái và bên dưới (giống tam giác Pascal):
\begin{figC}\begin{scope}[scale=0.75]
  \foreach \x/\y/\n in {0/0/1,1/0/1,2/0/1,3/0/1,4/0/1,0/1/1,1/1/2,2/1/3,3/1/4,0/2/1,1/2/3,2/2/6,
     0/3/1,1/3/4,0/4/1}{
    \fill[tmBlue!12] (\x,\y) circle (9pt);
    \node[font=\small\bfseries,text=Navy] at (\x,\y) {\n};}
  \draw[AnsRed,thick] (2,2) circle (10pt);
  \node[right=10pt,text=AnsRed,font=\small] at (2,2) {$R$};
  \node[right=10pt,font=\small] at (4,0) {$T$};\node[right=10pt,font=\small] at (3,1) {$S$};
  \node[right=10pt,font=\small] at (1,3) {$Q$};\node[right=10pt,font=\small] at (0,4) {$P$};
\end{scope}\end{figC}
Trên đường chéo cuối: $1+4+6+4+1=16$ (đúng bằng tổng số dãy) và tại $R$ có số $6$.
\begin{warnbox}Năm điểm $P,Q,R,S,T$ \emph{không} đồng khả năng! Không được trả lời ``$1$ trong $5$ điểm nên
xác suất là $\frac15$''. Điểm ở giữa ($R$) có nhiều đường đi tới nhất; hai điểm ở đầu ($P$ và $T$) mỗi điểm chỉ có một đường.\end{warnbox}
\verify{Xác suất dừng ở các điểm $P,Q,R,S,T$ lần lượt là $\frac1{16},\frac4{16},\frac6{16},\frac4{16},\frac1{16}$;
tổng bằng $\frac{16}{16}=1$, hợp lý.}
Giá trị $\frac38=\frac6{16}$ là phương án \textbf{(B)}.
\ans{(B) $\dfrac38$}
}

\begin{probbox}{17}{\relax}
Người ta dùng các que tăm giống hệt nhau để xếp thành các lưới ô vuông như hình dưới. Chẳng hạn, lưới
$1\times10$ (gồm $1$ hàng, $10$ cột ô vuông) cần tất cả $31$ que tăm. Hỏi để xếp lưới $43\times10$
($43$ hàng, mỗi hàng $10$ ô) thì cần bao nhiêu que tăm?
\begin{figC}
  \begin{scope}[shift={(0,2.4)}]\cpTooth{1}{10}{0.42}\node[right] at (4.4,0.21) {$1\times10$};\end{scope}
  \begin{scope}[shift={(0,1.05)}]\cpTooth{2}{10}{0.42}\node[right] at (4.4,0.42) {$2\times10$};\end{scope}
  \begin{scope}[shift={(0,-0.6)}]\cpTooth{3}{10}{0.42}\node[right] at (4.4,0.63) {$3\times10$};\end{scope}
\end{figC}
\choices{$913$}{$860$}{$871$}{$903$}{$946$}
\end{probbox}
\probsol{
\textbf{Phân tích đề.} Lưới $43\times10$ có $43$ hàng ô vuông, mỗi hàng $10$ ô. Các que tăm nằm theo hai
hướng: \emph{nằm ngang} và \emph{thẳng đứng}. Ta đếm riêng từng loại rồi cộng lại (quy tắc cộng).

\step{1}{Kiểm tra cách đếm trên lưới $1\times10$}
\emph{Que nằm ngang:} lưới có $2$ đường kẻ ngang (cạnh trên và cạnh dưới), mỗi đường gồm $10$ que:
$2\times10=20$ que.
\emph{Que thẳng đứng:} có $11$ đường kẻ dọc ($10$ ô cần $11$ vách ngăn), mỗi đường gồm $1$ que:
$11\times1=11$ que. Tổng: $20+11=31$ que, đúng như đề cho.
\begin{recallbox}Một hàng có $10$ ô thì cần $10+1=11$ vách dọc (giống bài toán trồng cây: $10$ khoảng cách
thì có $11$ cây). Tương tự, $43$ hàng ô thì có $43+1=44$ đường kẻ ngang.\end{recallbox}

\step{2}{Đếm que nằm ngang của lưới $43\times10$}
Có $43+1=44$ đường kẻ ngang, mỗi đường dài $10$ que: $44\times10=440$ que.

\step{3}{Đếm que thẳng đứng của lưới $43\times10$}
Có $10+1=11$ đường kẻ dọc, mỗi đường dài $43$ que (mỗi hàng ô góp $1$ que): $11\times43=473$ que.

\step{4}{Cộng lại}
$440+473=913$ que tăm.

\step{5}{Cách thứ hai: thêm từng hàng}
Lưới $1\times10$ có $31$ que. Mỗi khi đặt thêm một hàng mới lên trên, ta cần thêm $10$ que nằm ngang
(cạnh trên của hàng mới) và $11$ que thẳng đứng, tức $21$ que (cạnh dưới đã có sẵn). Từ $1$ hàng lên
$43$ hàng phải thêm $42$ hàng:
\[31+42\times21=31+882=913.\]
\begin{warnbox}Không được lấy $43\times31=1333$: khi chồng các hàng lên nhau, hai hàng kề nhau \emph{dùng chung}
một dãy $10$ que nằm ngang, nên nhân như vậy là đếm trùng.\end{warnbox}
\verify{Theo công thức $(\text{số hàng}+1)\times10+11\times\text{số hàng}$: lưới $2\times10$ cần
$3\times10+11\times2=52$ que, lưới $3\times10$ cần $4\times10+11\times3=73$ que; đếm trực tiếp trên hình cũng
được $52$ và $73$. Hai cách tính cho cùng kết quả $913$.}
Vậy chọn \textbf{(A)}.
\ans{(A) $913$ que tăm}
}

\begin{probbox}{18}{\relax}
Ở căng-tin trường bạn Braydon, món thịt có hai loại là thịt bò và thịt gà; trái cây có ba loại là táo, lê
và chuối. Braydon được phát ngẫu nhiên một suất ăn gồm một món thịt và một loại trái cây. Xác suất để
suất ăn của Braydon có chuối là
\choices{$\dfrac13$}{$\dfrac23$}{$\dfrac12$}{$\dfrac15$}{$\dfrac35$}
\end{probbox}
\probsol{
\textbf{Phân tích đề.} Một suất ăn $=$ (một món thịt, một loại trái cây). Suất ăn được chọn ngẫu nhiên
nên mọi suất ăn có thể đều đồng khả năng. Cần tìm xác suất suất ăn có chuối.

\step{1}{Liệt kê tất cả các suất ăn có thể}
Theo quy tắc nhân có $2\times3=6$ suất ăn. Ta lập bảng (hàng: thịt, cột: trái cây):
\begin{center}\begin{tabular}{|c|c|c|c|}\hline
 & Táo & Lê & \cellcolor{tmOrange!20}Chuối\\\hline
Bò & Bò--Táo & Bò--Lê & \cellcolor{tmOrange!20}Bò--Chuối\\\hline
Gà & Gà--Táo & Gà--Lê & \cellcolor{tmOrange!20}Gà--Chuối\\\hline
\end{tabular}\end{center}

\step{2}{Đếm các suất ăn có chuối}
Cột ``Chuối'' có $2$ suất: Bò--Chuối và Gà--Chuối.

\step{3}{Tính xác suất}
\[P(\text{có chuối})=\frac{2}{6}=\frac13.\]
\begin{recallbox}Rút gọn phân số: chia cả tử và mẫu cho ước chung lớn nhất: $\dfrac26=\dfrac{2:2}{6:2}=\dfrac13$.\end{recallbox}
\begin{tipbox}Món thịt không ảnh hưởng gì tới trái cây. Có thể nghĩ gọn: trái cây được chọn là một trong
$3$ loại đồng khả năng, chuối là $1$ trong $3$ loại, nên xác suất là $\frac13$.\end{tipbox}
\begin{warnbox}Đừng cộng ``$2$ loại thịt $+\,3$ loại trái cây $=5$'' rồi trả lời $\frac15$: số suất ăn là
$2\times3=6$ chứ không phải $2+3=5$.\end{warnbox}
\verify{Xác suất có táo, có lê, có chuối đều bằng $\frac13$; tổng $\frac13+\frac13+\frac13=1$.}
Vậy chọn \textbf{(A)}.
\ans{(A) $\dfrac13$}
}

\begin{probbox}{19}{\relax}
Người ta làm một con xúc xắc đặc biệt có sáu mặt, mỗi mặt ghi một số (các số có thể trùng nhau). Khi gieo
xúc xắc này, xác suất được một số chia hết cho $3$ là $\frac12$, còn xác suất được một số chẵn là $\frac13$.
Sáu số trên các mặt có thể là
\choicesII{$1,2,3,5,5,6$}{$1,2,3,3,5,6$}{$1,2,3,4,6,6$}{$1,2,3,3,4,6$}{$2,3,3,3,5,6$}
\end{probbox}
\probsol{
\textbf{Phân tích đề.} Xúc xắc có $6$ mặt đồng khả năng. Từ hai xác suất đã cho, ta suy ra \emph{số mặt}
ghi số chia hết cho $3$ và số mặt ghi số chẵn, rồi kiểm tra từng phương án.

\step{1}{Đổi xác suất thành số mặt}
\begin{recallbox}Với $6$ mặt đồng khả năng: $P=\dfrac{\text{số mặt thuận lợi}}{6}$, nên
số mặt thuận lợi $=P\times6$.\end{recallbox}
Số mặt chia hết cho $3$: $\frac12\times6=3$ mặt.\quad Số mặt chẵn: $\frac13\times6=2$ mặt.

Vậy ta cần: \textbf{đúng $3$ mặt chia hết cho $3$} và \textbf{đúng $2$ mặt chẵn}.

\step{2}{Kiểm tra từng phương án}
Trong các số từ $1$ đến $6$: các số chia hết cho $3$ là $3$ và $6$; các số chẵn là $2$, $4$, $6$.
Chú ý số $6$ vừa chẵn vừa chia hết cho $3$, nên phải đếm nó ở \emph{cả hai} cột.
\begin{center}\begin{tabular}{|c|c|c|c|c|}\hline
Phương án & Các mặt & Mặt chia hết cho $3$ & Mặt chẵn & Đạt?\\\hline
(A) & $1,2,3,5,5,6$ & $3,6$ \;($2$ mặt) & $2,6$ \;($2$ mặt) & Không\\\hline
\rowcolor{tmGreen!12}(B) & $1,2,3,3,5,6$ & $3,3,6$ \;($3$ mặt) & $2,6$ \;($2$ mặt) & Có\\\hline
(C) & $1,2,3,4,6,6$ & $3,6,6$ \;($3$ mặt) & $2,4,6,6$ \;($4$ mặt) & Không\\\hline
(D) & $1,2,3,3,4,6$ & $3,3,6$ \;($3$ mặt) & $2,4,6$ \;($3$ mặt) & Không\\\hline
(E) & $2,3,3,3,5,6$ & $3,3,3,6$ \;($4$ mặt) & $2,6$ \;($2$ mặt) & Không\\\hline
\end{tabular}\end{center}
\begin{warnbox}Nếu quên rằng $6$ là số chẵn, ta sẽ đếm sai số mặt chẵn của (D) thành $2$ và chọn nhầm (D).
Một số có thể thuộc nhiều nhóm cùng lúc!\end{warnbox}
\verify{Với (B): $P(\text{chia hết cho }3)=\frac36=\frac12$ và $P(\text{chẵn})=\frac26=\frac13$, đúng cả hai điều kiện.}
Chỉ có phương án \textbf{(B)} thoả mãn.
\ans{(B) $1,2,3,3,5,6$}
}

\begin{probbox}{20}{\stars{1}}
Ba tuyến xe buýt chạy trên các con đường như hình vẽ, cùng đi qua bến $X$ (các đoạn có cùng số vạch
nhỏ thì dài bằng nhau: $PX=XS$, $QX=XT$, $RX=XU$).
\begin{itemize}[leftmargin=1.2em,itemsep=1pt,topsep=2pt]
\item Xe $A$ chạy theo lộ trình $P\to X\to S\to X\to P$, mỗi vòng mất $12$ phút.
\item Xe $B$ chạy theo lộ trình $Q\to X\to T\to X\to Q$, mỗi vòng mất $20$ phút.
\item Xe $C$ chạy theo lộ trình $R\to X\to U\to X\to R$, mỗi vòng mất $28$ phút.
\end{itemize}
Lúc $13$\,h, ba xe $A$, $B$, $C$ lần lượt xuất phát từ $P$, $Q$, $R$. Mỗi xe chạy với vận tốc không đổi,
quay đầu ngay lập tức ở hai đầu lộ trình và chạy liên tục đến $23$\,h. Hỏi trong khoảng thời gian từ
$17$\,h đến $22$\,h, có bao nhiêu thời điểm mà có từ hai xe trở lên cùng có mặt tại bến $X$?
\begin{figC}\begin{scope}[scale=0.9]
  \draw[Ink,thick] (0,1.6)--(0,-1.6);\draw[Ink,thick] (-1.9,0)--(1.9,0);\draw[Ink,thick] (-2,-1.35)--(2,1.35);
  \foreach \p in {(0,1.6),(0,-1.6),(-1.9,0),(1.9,0),(-2,-1.35),(2,1.35),(0,0)}{\fill[Navy] \p circle (2pt);}
  \node[above] at (0,1.6) {$P$};\node[below] at (0,-1.6) {$S$};\node[left] at (-1.9,0) {$Q$};
  \node[right] at (1.9,0) {$T$};\node[below left] at (-2,-1.35) {$R$};\node[above right] at (2,1.35) {$U$};
  \node[below right] at (0.05,-0.02) {$X$};
  % ticks
  \draw[tmBlue,thick] (-0.12,0.8)--(0.12,0.8);\draw[tmBlue,thick] (-0.12,-0.8)--(0.12,-0.8);
  \foreach \s in {-1,1}{\draw[tmOrange,thick] (\s*0.95-0.04,-0.12)--(\s*0.95-0.04,0.12);
                        \draw[tmOrange,thick] (\s*0.95+0.04,-0.12)--(\s*0.95+0.04,0.12);}
  \foreach \s in {-1,1}{\foreach \d in {-0.09,0,0.09}{
     \draw[tmGreen!80!black,thick] ($(\s*1,\s*0.675)+(\d*0.829,\d*0.559)+(0.067,-0.1)$)--++(-0.134,0.2);}}
\end{scope}\end{figC}
\choices{$18$}{$19$}{$20$}{$21$}{$22$}
\end{probbox}
\probsol{
\textbf{Phân tích đề.} Ta đo thời gian bằng \emph{số phút kể từ $13$\,h} và gọi số đó là $t$. Cần tìm
những giá trị $t$ mà có ít nhất hai xe cùng ở $X$, với $t$ nằm trong khoảng từ $17$\,h đến $22$\,h.

\step{1}{Đổi khoảng thời gian cần xét ra số phút}
$17$\,h là $4$ giờ sau $13$\,h: $t=4\times60=240$. \quad $22$\,h là $9$ giờ sau $13$\,h: $t=9\times60=540$.
Vậy ta xét $240\le t\le540$.

\step{2}{Xe $A$ có mặt ở $X$ vào những lúc nào?}
Một vòng của xe $A$ gồm $4$ đoạn bằng nhau: $P\to X$, $X\to S$, $S\to X$, $X\to P$ (vì $PX=XS$).
Xe chạy đều nên mỗi đoạn mất $12:4=3$ phút. Xe $A$ tới $X$ ở cuối đoạn thứ nhất và cuối đoạn thứ ba
của mỗi vòng, tức là lúc
\[t=3,\ 9,\ 15,\ 21,\ 27,\ \dots\]
Đó chính là các \textbf{bội lẻ của $3$} ($3\times1$, $3\times3$, $3\times5$, \dots). Ở cuối đoạn thứ hai và thứ tư
xe ở $S$ hoặc $P$, không ở $X$.

\step{3}{Tương tự với xe $B$ và xe $C$}
Xe $B$: mỗi đoạn mất $20:4=5$ phút, xe ở $X$ lúc $t=5,15,25,35,\dots$ -- các \textbf{bội lẻ của $5$}.\\
Xe $C$: mỗi đoạn mất $28:4=7$ phút, xe ở $X$ lúc $t=7,21,35,49,\dots$ -- các \textbf{bội lẻ của $7$}.

\step{4}{Khi nào hai xe gặp nhau ở $X$?}
\begin{recallbox}Một số chia hết cho cả $3$ và $5$ thì chia hết cho $\text{BCNN}(3,5)=15$. Tương tự:
chia hết cho cả $3$ và $7$ $\Leftrightarrow$ chia hết cho $21$; chia hết cho cả $5$ và $7$ $\Leftrightarrow$
chia hết cho $35$.\end{recallbox}
Xe $A$ và xe $B$ cùng ở $X$ khi $t$ vừa là bội lẻ của $3$ vừa là bội lẻ của $5$, tức $t$ là \textbf{số lẻ chia hết cho $15$}
(bội lẻ của $15$). Tương tự: $A$ và $C$ gặp nhau khi $t$ là bội lẻ của $21$; $B$ và $C$ gặp nhau khi $t$ là bội lẻ của $35$.
\begin{proofbox}Bội lẻ của $15$ là số lẻ chia hết cho $15$, nên nó vừa là số lẻ chia hết cho $3$ (bội lẻ của $3$), vừa là
số lẻ chia hết cho $5$ (bội lẻ của $5$). Ngược lại, một số vừa là bội lẻ của $3$ vừa là bội lẻ của $5$ thì lẻ và chia
hết cho $15$.\end{proofbox}

\step{5}{Liệt kê trong khoảng $240\le t\le540$}
\begin{center}\begin{tabular}{|l|l|c|}\hline
Hai xe & Các giá trị $t$ (bội lẻ) & Số lần\\\hline
$A$ và $B$ ($15$) & $255, 285, 315, 345, 375, 405, 435, 465, 495, 525$ & $10$\\\hline
$A$ và $C$ ($21$) & $273, 315, 357, 399, 441, 483, 525$ & $7$\\\hline
$B$ và $C$ ($35$) & $245, 315, 385, 455, 525$ & $5$\\\hline
\end{tabular}\end{center}
(Ví dụ hàng đầu: $15\times17=255$ là bội lẻ đầu tiên lớn hơn $240$; sau đó cứ cộng thêm $30$, vì giữa hai
bội lẻ liên tiếp của $15$ cách nhau $2\times15=30$.)

\step{6}{Bỏ các thời điểm bị đếm trùng}
Cộng ba hàng được $10+7+5=22$, nhưng các thời điểm $315$ và $525$ xuất hiện ở \emph{cả ba hàng}
(đó là bội lẻ của $3\times5\times7=105$: lúc ấy cả ba xe cùng ở $X$). Mỗi thời điểm này bị đếm $3$ lần
trong khi chỉ được đếm $1$ lần, nên phải bớt $2+2=4$:
\[22-4=18 \text{ thời điểm.}\]
\begin{warnbox}Đề hỏi số \emph{thời điểm}, không phải số ``cặp xe gặp nhau''. Lúc ba xe cùng ở $X$ chỉ là
\emph{một} thời điểm. Cộng thẳng $10+7+5=22$ rồi chọn (E) là sai.\end{warnbox}
\verify{Đổi ra giờ: $17$\,h\,05; $17$\,h\,15; $17$\,h\,33; $17$\,h\,45; $18$\,h\,15; $18$\,h\,45; $18$\,h\,57;
$19$\,h\,15; $19$\,h\,25; $19$\,h\,39; $19$\,h\,45; $20$\,h\,15; $20$\,h\,21; $20$\,h\,35; $20$\,h\,45; $21$\,h\,03;
$21$\,h\,15; $21$\,h\,45 -- đếm được đúng $18$ thời điểm, tất cả đều nằm giữa $17$\,h và $22$\,h.}
Vậy chọn \textbf{(A)}.
\ans{(A) $18$}
}

\begin{probbox}{21}{\relax}
Mỗi mặt của một hình lập phương được sơn đúng một màu. Cần ít nhất bao nhiêu màu để sơn hình lập
phương sao cho hai mặt bất kì có chung một cạnh thì có màu khác nhau?
\choices{$2$}{$3$}{$4$}{$5$}{$6$}
\end{probbox}
\probsol{
\textbf{Phân tích đề.} Hai mặt ``có chung một cạnh'' là hai mặt kề nhau. Hai mặt đối diện (như mặt trên và
mặt dưới) không có cạnh chung nên được phép cùng màu. Ta cần (1) chỉ ra một cách tô dùng ít màu, và
(2) chứng minh không thể dùng ít màu hơn.

\step{1}{Hai màu là không đủ}
\begin{proofbox}Xét một đỉnh bất kì của hình lập phương, chẳng hạn đỉnh trên--trước--phải. Có ba mặt gặp nhau ở
đỉnh đó: mặt trên, mặt trước, mặt phải. Hai mặt bất kì trong ba mặt này đều có chung một cạnh, nên ba mặt
này phải có ba màu đôi một khác nhau. Do đó cần ít nhất $3$ màu.\end{proofbox}

\step{2}{Ba màu là đủ}
Tô \emph{hai mặt đối diện cùng một màu}: mặt trên và mặt dưới màu xanh dương, mặt trước và mặt sau màu
cam, mặt trái và mặt phải màu xanh lá. Mỗi mặt kề với $4$ mặt, đó chính là bốn mặt \emph{không} đối diện
với nó, và bốn mặt này mang hai màu kia. Vậy không có hai mặt kề nhau nào cùng màu.
\begin{figC}
  \begin{scope}
   \fill[tmOrange!55,draw=Ink,thick] (0,0)--(1.6,0)--(1.6,1.6)--(0,1.6)--cycle;
   \fill[tmBlue!55,draw=Ink,thick] (0,1.6)--(1.6,1.6)--(2.4,2.3)--(0.8,2.3)--cycle;
   \fill[tmGreen!55,draw=Ink,thick] (1.6,0)--(2.4,0.7)--(2.4,2.3)--(1.6,1.6)--cycle;
   \node at (0.8,0.8) {trước};\node at (1.2,1.95) {trên};\node[rotate=0] at (2,1.15) {phải};
  \end{scope}
  \begin{scope}[shift={(4,0.1)},scale=0.5]
   \foreach \x/\y/\c in {1/2/tmBlue,0/1/tmGreen,1/1/tmOrange,2/1/tmGreen,3/1/tmOrange,1/0/tmBlue}{
     \fill[\c!55,draw=Ink,thick] (\x,\y) rectangle ++(1,1);}
   \node[font=\scriptsize] at (2,-0.45) {hình khai triển};
  \end{scope}
\end{figC}

\step{3}{Kết luận}
Cần ít nhất $3$ màu (bước 1) và $3$ màu là đủ (bước 2), nên số màu nhỏ nhất là $3$.
\begin{warnbox}Chỉ tìm được một cách tô bằng $3$ màu thì \emph{chưa đủ} để khẳng định $3$ là nhỏ nhất; phải
giải thích vì sao $2$ màu không được (lập luận ba mặt chung đỉnh). Ngược lại, nghĩ rằng ``$6$ mặt phải có
$6$ màu'' là quên rằng mặt đối diện được phép cùng màu.\end{warnbox}
\verify{Trên hình khai triển, mỗi ô vuông chỉ chung cạnh với các ô khác màu (khi gấp lại, hai ô cùng màu
là hai mặt đối diện).}
Vậy chọn \textbf{(B)}.
\ans{(B) $3$ màu}
}

\begin{probbox}{22}{\relax}
Có bao nhiêu cặp số $(m,n)$ khác nhau lấy từ các số nguyên $1,2,3,\dots,20$ sao cho $m<n$ và $m+n$ là số chẵn?
\choices{$55$}{$90$}{$140$}{$110$}{$50$}
\end{probbox}
\probsol{
\textbf{Phân tích đề.} Ta chọn hai số khác nhau trong $20$ số; vì yêu cầu $m<n$ nên mỗi cặp hai số chỉ được
đếm \emph{một lần} (số nhỏ là $m$, số lớn là $n$). Điều kiện: tổng $m+n$ chẵn.

\step{1}{Khi nào tổng hai số là số chẵn?}
\begin{recallbox}chẵn $+$ chẵn $=$ chẵn;\quad lẻ $+$ lẻ $=$ chẵn;\quad chẵn $+$ lẻ $=$ lẻ.\end{recallbox}
Vậy $m+n$ chẵn khi và chỉ khi $m$ và $n$ \textbf{cùng tính chẵn lẻ} (cùng chẵn hoặc cùng lẻ).

\step{2}{Chia $20$ số thành hai nhóm}
Nhóm chẵn: $2,4,6,\dots,20$ -- có $10$ số.\quad Nhóm lẻ: $1,3,5,\dots,19$ -- có $10$ số.

\step{3}{Đếm số cặp trong nhóm chẵn}
Chọn số thứ nhất có $10$ cách, số thứ hai (khác số thứ nhất) có $9$ cách: $10\times9=90$ cách chọn \emph{có thứ tự}.
Nhưng mỗi cặp, ví dụ $\{2,4\}$, bị đếm hai lần ($2$ rồi $4$, hoặc $4$ rồi $2$), trong khi điều kiện $m<n$
chỉ cho phép một thứ tự $(2,4)$. Vậy số cặp là $\dfrac{10\times9}{2}=45$.
\begin{recallbox}Số cách chọn $2$ phần tử (không kể thứ tự) từ $k$ phần tử là $\dfrac{k\times(k-1)}{2}$.\end{recallbox}

\step{4}{Đếm số cặp trong nhóm lẻ và cộng lại}
Tương tự, nhóm lẻ cho $45$ cặp. Một cặp không thể vừa ở nhóm chẵn vừa ở nhóm lẻ, nên theo quy tắc cộng:
$45+45=90$ cặp.
\begin{warnbox}Quên chia cho $2$ sẽ được $90+90=180$ (không có trong đáp án); quên mất nhóm lẻ sẽ được
$45$. Hãy luôn tự hỏi: ``Mình có đếm một cặp hai lần không? Mình có bỏ sót trường hợp nào không?''\end{warnbox}
\verify{Đếm bằng phần bù: tổng số cặp $m<n$ là $\dfrac{20\times19}{2}=190$. Số cặp khác tính chẵn lẻ (một chẵn, một lẻ) là
$10\times10=100$. Vậy số cặp cùng tính chẵn lẻ là $190-100=90$. Khớp!}
Vậy chọn \textbf{(B)}.
\ans{(B) $90$ cặp}
}

\begin{probbox}{23}{\stars{1}}
Hai số nguyên dương có hai chữ số, khác nhau, được gọi là một \emph{cặp số đảo} nếu số này nhận được từ số kia bằng
cách đổi chỗ hai chữ số. Ví dụ: $52$ và $25$ là một cặp số đảo. Số $2015$ có một tính chất đặc biệt: nó bằng
tích của ba số nguyên tố khác nhau, trong đó có hai số tạo thành một cặp số đảo ($2015=5\times13\times31$).
Kể cả $2015$, có bao nhiêu số nguyên dương nhỏ hơn $10\,000$ có cùng tính chất này?
\choices{$18$}{$14$}{$20$}{$17$}{$19$}
\end{probbox}
\probsol{
\textbf{Phân tích đề.} Ta cần đếm các số $N<10\,000$ có dạng $N=p\times q\times r$, với $p,q,r$ là ba số nguyên
tố khác nhau, trong đó hai số (gọi là $p$ và $q$) là một cặp số đảo có hai chữ số. Kế hoạch: tìm tất cả các
cặp số đảo gồm hai số nguyên tố, rồi với mỗi cặp tìm các số nguyên tố thứ ba phù hợp.

\step{1}{Tìm các cặp số đảo mà cả hai số đều là số nguyên tố}
\begin{recallbox}Số nguyên tố là số tự nhiên lớn hơn $1$ chỉ có hai ước là $1$ và chính nó. Các số nguyên tố có
hai chữ số: $11, 13, 17, 19, 23, 29, 31, 37, 41, 43, 47, 53, 59, 61, 67, 71, 73, 79, 83, 89, 97$.\end{recallbox}
\begin{tipbox}Số đảo của $p$ có chữ số hàng đơn vị chính là chữ số hàng chục của $p$. Số nguyên tố có hai chữ số
phải tận cùng bằng $1,3,7$ hoặc $9$, nên chỉ cần xét các số nguyên tố bắt đầu bằng $1,3,7,9$:
$11,13,17,19,31,37,71,73,79,97$.\end{tipbox}
Kiểm tra số đảo của từng số:
\begin{center}\begin{tabular}{|c|c|c|c|c|c|c|c|c|c|c|}\hline
$p$ & $11$ & $13$ & $17$ & $19$ & $31$ & $37$ & $71$ & $73$ & $79$ & $97$\\\hline
số đảo & $11$ & $31$ & $71$ & $91$ & $13$ & $73$ & $17$ & $37$ & $97$ & $79$\\\hline
nguyên tố? & trùng & có & có & không & có & có & có & có & có & có\\\hline
\end{tabular}\end{center}
($91=7\times13$ không là số nguyên tố; $11$ đảo lại vẫn là $11$, không phải hai số \emph{khác nhau}.)
Vậy có đúng $4$ cặp: $\{13,31\}$, $\{17,71\}$, $\{37,73\}$, $\{79,97\}$.

\step{2}{Tính tích của mỗi cặp}
$13\times31=403$;\quad $17\times71=1207$;\quad $37\times73=2701$;\quad $79\times97=7663$.

\step{3}{Tìm số nguyên tố thứ ba $r$ cho từng cặp}
Cần $r$ là số nguyên tố, khác hai số của cặp, và tích $<10\,000$.
\begin{itemize}[leftmargin=1.2em,itemsep=1pt,topsep=2pt]
\item Cặp $\{13,31\}$: cần $403\times r<10\,000$. Vì $403\times24=9672$ và $403\times25=10\,075$ nên $r\le23$.
Các số nguyên tố $\le23$ là $2,3,5,7,11,13,17,19,23$; bỏ $13$ (đã có trong cặp): còn $8$ số, cho
$806;\ 1209;\ 2015;\ 2821;\ 4433;\ 6851;\ 7657;\ 9269$.
\item Cặp $\{17,71\}$: cần $1207\times r<10\,000$, mà $1207\times8=9656$, $1207\times9=10\,863$ nên $r\le8$:
$r\in\{2,3,5,7\}$, cho $2414;\ 3621;\ 6035;\ 8449$ ($4$ số).
\item Cặp $\{37,73\}$: $2701\times3=8103<10\,000$, $2701\times4>10\,000$ nên $r\in\{2,3\}$: $5402;\ 8103$ ($2$ số).
\item Cặp $\{79,97\}$: ngay cả $7663\times2=15\,326>10\,000$, nên không có số nào.
\end{itemize}

\step{4}{Kiểm tra không có số nào bị đếm hai lần và cộng lại}
\begin{proofbox}Mỗi số chỉ có một cách phân tích ra thừa số nguyên tố. Nếu một số xuất hiện ở hai danh sách thì
trong ba thừa số nguyên tố của nó có hai cặp số đảo khác nhau; hai cặp này phải dùng chung một số nguyên tố.
Nhưng nhìn bảng ở bước 1, mỗi số nguyên tố chỉ thuộc đúng một cặp. Vậy không có số nào bị đếm hai lần.\end{proofbox}
Tổng cộng: $8+4+2+0=14$ số.
\begin{warnbox}Hai lỗi hay gặp: (1) quên rằng số thứ ba phải \emph{khác} hai số của cặp, nên đếm thêm
$403\times13=5239=13\times13\times31$ (sai, vì $13$ lặp lại); (2) quên số nguyên tố $2$.\end{warnbox}
\verify{$2015=403\times5$ có trong danh sách. Kiểm tra vài số: $9269=13\times31\times23$, $8449=17\times71\times7$,
$8103=37\times73\times3$; tất cả đều nhỏ hơn $10\,000$.}
Vậy chọn \textbf{(B)}.
\ans{(B) $14$}
}

\begin{probbox}{24}{\relax}
Hình bên được ghép từ $8$ hình bình hành nhỏ giống hệt nhau. Kể cả $8$ hình bình hành nhỏ đó, trong
hình có tất cả bao nhiêu hình bình hành?
\begin{figC}
  \foreach \j in {0,1,2}{\draw[Ink,thick] (\j*0.55,\j*0.8)--(4.4+\j*0.55,\j*0.8);}
  \foreach \i in {0,...,4}{\draw[Ink,thick] (\i*1.1,0)--(\i*1.1+1.1,1.6);}
\end{figC}
\choices{$29$}{$30$}{$26$}{$27$}{$28$}
\end{probbox}
\probsol{
\textbf{Phân tích đề.} Hình gồm $2$ hàng, mỗi hàng $4$ hình bình hành nhỏ. Mỗi hình bình hành trong hình được
tạo bởi các đoạn thẳng có sẵn: hai cạnh nằm trên \emph{đường nằm ngang}, hai cạnh nằm trên \emph{đường xiên}.

\step{1}{Đếm các đường thẳng của hình}
Đánh số: có $3$ đường nằm ngang (đáy, giữa, trên) và $5$ đường xiên (song song với nhau).
\begin{figC}
  \fill[tmGreen!30] (1.1,0)--(3.3,0)--(3.85,0.8)--(1.65,0.8)--cycle;
  \foreach \j/\l in {0/1,1/2,2/3}{\draw[tmBlue,thick] (\j*0.55,\j*0.8)--(4.4+\j*0.55,\j*0.8);
     \node[left,text=tmBlue,font=\scriptsize] at (\j*0.55,\j*0.8) {ngang \l};}
  \foreach \i in {0,...,4}{\draw[tmOrange,thick] (\i*1.1,0)--(\i*1.1+1.1,1.6);
     \pgfmathtruncatemacro\k{\i+1}\node[below,text=tmOrange!80!black,font=\scriptsize] at (\i*1.1,0) {x\k};}
\end{figC}
Ví dụ: hình bình hành tô xanh lá được xác định bởi đường ngang $1$, ngang $2$ và đường xiên x$2$, x$4$.

\step{2}{Mỗi hình bình hành ứng với một cách chọn $2$ đường ngang và $2$ đường xiên}
Chọn $2$ trong $3$ đường ngang làm cạnh dưới và cạnh trên, chọn $2$ trong $5$ đường xiên làm cạnh trái và cạnh phải,
ta được đúng một hình bình hành; ngược lại mỗi hình bình hành trong hình cho đúng một cách chọn như vậy.
\begin{recallbox}Số cách chọn $2$ đường trong $k$ đường là $\dfrac{k\times(k-1)}{2}$.\end{recallbox}
Chọn $2$ đường ngang: $\dfrac{3\times2}{2}=3$ cách. \quad Chọn $2$ đường xiên: $\dfrac{5\times4}{2}=10$ cách.

\step{3}{Áp dụng quy tắc nhân}
Số hình bình hành: $3\times10=30$.

\step{4}{Kiểm tra bằng cách đếm theo kích thước}
Gọi hình bình hành ``cao $a$, rộng $b$'' là hình gồm $a$ hàng và $b$ cột hình nhỏ:
\begin{center}\begin{tabular}{|c|c|c|c|c|c|}\hline
 & rộng $1$ & rộng $2$ & rộng $3$ & rộng $4$ & Cộng\\\hline
cao $1$ (có $2$ vị trí theo hàng) & $2\times4=8$ & $2\times3=6$ & $2\times2=4$ & $2\times1=2$ & $20$\\\hline
cao $2$ (có $1$ vị trí theo hàng) & $4$ & $3$ & $2$ & $1$ & $10$\\\hline
\end{tabular}\end{center}
\begin{warnbox}Đừng chỉ đếm $8$ hình nhỏ cộng vài hình lớn ``dễ thấy''. Hãy đếm có hệ thống: theo kích thước
(bảng trên) hoặc theo cách chọn các đường thẳng.\end{warnbox}
\verify{Tổng theo bảng: $20+10=30$, trùng với cách tính $3\times10$.}
Vậy chọn \textbf{(B)}.
\ans{(B) $30$ hình bình hành}
}

\begin{probbox}{25}{\stars{1}}
Trong một cái lọ có $2$ viên bi đỏ và $2$ viên bi xanh, ngoài ra không có viên bi nào khác. Bạn Angie lấy ngẫu
nhiên cùng lúc $2$ viên bi từ lọ. Nếu hai viên cùng màu, Angie bỏ ra ngoài một viên và trả viên kia vào lọ.
Nếu hai viên khác màu, Angie bỏ viên bi đỏ ra ngoài và trả viên bi xanh vào lọ. Angie làm như vậy tất cả ba lần.
Xác suất để viên bi cuối cùng còn lại trong lọ là bi đỏ bằng
\choices{$\dfrac12$}{$\dfrac14$}{$\dfrac23$}{$\dfrac13$}{$0$}
\end{probbox}
\probsol{
\textbf{Phân tích đề.} Mỗi lần lấy, trong lọ bớt đi đúng $1$ viên bi. Sau $3$ lần, từ $4$ viên còn lại $1$ viên.
Ta theo dõi \emph{trạng thái} của lọ, viết (số bi đỏ, số bi xanh), bắt đầu từ $(2\text{Đ},2\text{X})$.

\step{1}{Lần lấy thứ nhất, từ $(2\text{Đ},2\text{X})$}
Gọi các bi là Đ$_1$, Đ$_2$, X$_1$, X$_2$. Có $6$ cách lấy $2$ viên, đồng khả năng:
Đ$_1$Đ$_2$;\ X$_1$X$_2$;\ Đ$_1$X$_1$;\ Đ$_1$X$_2$;\ Đ$_2$X$_1$;\ Đ$_2$X$_2$.
\begin{itemize}[leftmargin=1.2em,itemsep=1pt,topsep=2pt]
\item Hai bi đỏ ($1$ cách): bỏ $1$ bi đỏ $\to(1\text{Đ},2\text{X})$.
\item Hai bi xanh ($1$ cách): bỏ $1$ bi xanh $\to(2\text{Đ},1\text{X})$.
\item Khác màu ($4$ cách): bỏ bi đỏ $\to(1\text{Đ},2\text{X})$.
\end{itemize}
Vậy sau lần $1$: trạng thái $(1\text{Đ},2\text{X})$ với xác suất $\frac56$, trạng thái $(2\text{Đ},1\text{X})$ với xác suất $\frac16$.

\step{2}{Lần lấy thứ hai}
Từ $(1\text{Đ},2\text{X})$: có $3$ cách lấy; $2$ cách khác màu $\to(0\text{Đ},2\text{X})$; $1$ cách hai bi xanh $\to(1\text{Đ},1\text{X})$.\\
Từ $(2\text{Đ},1\text{X})$: có $3$ cách lấy; $1$ cách hai bi đỏ $\to(1\text{Đ},1\text{X})$; $2$ cách khác màu $\to(1\text{Đ},1\text{X})$.
Vậy từ $(2\text{Đ},1\text{X})$ chắc chắn sang $(1\text{Đ},1\text{X})$.

\step{3}{Lần lấy thứ ba}
Từ $(0\text{Đ},2\text{X})$: lấy hai bi xanh, bỏ một $\to(0\text{Đ},1\text{X})$.\\
Từ $(1\text{Đ},1\text{X})$: lấy hai bi khác màu, bỏ bi đỏ $\to(0\text{Đ},1\text{X})$.
\begin{figC}
  \tikzset{st/.style={draw=tmBlue,fill=tmBlue!8,rounded corners=3pt,inner sep=3pt,font=\footnotesize}}
  \node[st] (a) at (0,0) {$2$Đ, $2$X};
  \node[st] (b) at (2.8,0.9) {$1$Đ, $2$X};\node[st] (c) at (2.8,-0.9) {$2$Đ, $1$X};
  \node[st] (d) at (5.6,0.9) {$0$Đ, $2$X};\node[st] (e) at (5.6,-0.9) {$1$Đ, $1$X};
  \node[st,draw=AnsRed,fill=AnsRed!8] (f) at (8.2,0) {$0$Đ, $1$X};
  \draw[->,thick,Slate] (a)--node[above,font=\scriptsize]{$\frac56$}(b);
  \draw[->,thick,Slate] (a)--node[below,font=\scriptsize]{$\frac16$}(c);
  \draw[->,thick,Slate] (b)--node[above,font=\scriptsize]{$\frac23$}(d);
  \draw[->,thick,Slate] (b)--node[above right,font=\scriptsize,pos=0.3]{$\frac13$}(e);
  \draw[->,thick,Slate] (c)--node[below,font=\scriptsize]{$1$}(e);
  \draw[->,thick,Slate] (d)--node[above,font=\scriptsize]{$1$}(f);
  \draw[->,thick,Slate] (e)--node[below,font=\scriptsize]{$1$}(f);
  \node[font=\scriptsize,text=Slate] at (1.4,-1.7) {lần 1};\node[font=\scriptsize,text=Slate] at (4.2,-1.7) {lần 2};
  \node[font=\scriptsize,text=Slate] at (6.9,-1.7) {lần 3};
\end{figC}
Mọi nhánh đều kết thúc ở $(0\text{Đ},1\text{X})$: viên bi cuối cùng \textbf{luôn} là bi xanh.

\step{4}{Giải thích vì sao -- không cần vẽ cây}
\begin{proofbox}Hãy nhìn vào số bi xanh. Mỗi khi trong hai viên lấy ra có ít nhất một bi xanh, Angie luôn trả lại vào lọ một
bi xanh (nếu hai bi xanh thì trả một viên; nếu khác màu thì trả viên xanh). Số bi xanh chỉ giảm khi lấy được hai bi xanh,
và khi đó vẫn còn lại ít nhất viên vừa trả vào. Vậy trong lọ \emph{luôn còn ít nhất một bi xanh}. Khi chỉ còn một viên, viên đó phải
là bi xanh.\end{proofbox}
Do đó biến cố ``viên còn lại là bi đỏ'' là \emph{biến cố không thể}, có xác suất bằng $0$.
\begin{recallbox}Biến cố không thể (không bao giờ xảy ra) có xác suất $0$; biến cố chắc chắn có xác suất $1$.\end{recallbox}
\begin{warnbox}Thấy ``$2$ đỏ, $2$ xanh'' rồi đoán ngay xác suất $\frac12$ là bẫy: luật bỏ bi \emph{không đối xứng}
giữa hai màu (khi khác màu thì luôn bỏ bi đỏ).\end{warnbox}
\verify{Tổng xác suất tới trạng thái cuối $(0\text{Đ},1\text{X})$: $\frac56\cdot\frac23+\frac56\cdot\frac13+\frac16\cdot1=\frac{10}{18}+\frac{5}{18}+\frac{3}{18}=1$.}
Vậy chọn \textbf{(E)}.
\ans{(E) $0$}
}

\begin{probbox}{26}{\stars{1}}
Trên lưới ô vuông $10\times10$ ở hình bên, cho tam giác $FGH$. Điểm $P$ có thể là bất kì điểm nào trong
$41$ giao điểm của các đường kẻ lưới nằm \emph{bên trong} (không nằm trên cạnh) tam giác $FGH$. Với mỗi vị trí của $P$,
ta có đúng ba tam giác: $FPG$, $GPH$, $HPF$. Trong tất cả $41\times3=123$ tam giác đó, có bao nhiêu tam giác có diện
tích bằng đúng một nửa diện tích tam giác $FGH$?
\begin{figC}\begin{scope}[scale=0.3]
  \draw[Line,thin] (0,0) grid (10,10);
  \draw[Ink,line width=1.3pt] (0,0)--(10,0)--(5,10)--cycle;
  \node[above,font=\scriptsize] at (5,10) {$F$};\node[below left,font=\scriptsize] at (0,0) {$G$};
  \node[below right,font=\scriptsize] at (10,0) {$H$};
\end{scope}\end{figC}
\choices{$5$}{$7$}{$3$}{$11$}{$9$}
\end{probbox}
\probsol{
\textbf{Phân tích đề.} Lấy cạnh mỗi ô vuông làm đơn vị và đặt $G$ ở góc dưới trái. Khi đó $G$, $H$ ở hai góc dưới
($GH=10$), $F$ ở chính giữa cạnh trên. Mỗi điểm lưới được gọi tên bằng hai số (số ô sang phải, số ô lên trên); ví dụ
$G=(0,0)$, $H=(10,0)$, $F=(5,10)$.

\step{1}{Tính diện tích tam giác $FGH$ và một nửa của nó}
\begin{recallbox}Diện tích tam giác $=\dfrac12\times\text{cạnh đáy}\times\text{chiều cao tương ứng}$.\end{recallbox}
$S_{FGH}=\frac12\times10\times10=50$. Một nửa là $25$.

\step{2}{Tam giác $GPH$: dễ nhất}
Tam giác $GPH$ có đáy $GH=10$ và chiều cao là khoảng cách từ $P$ xuống $GH$, tức là ``độ cao'' $h$ của $P$
(số ô tính từ đáy). $S_{GPH}=\frac12\times10\times h=5h$. Ta cần $5h=25$, tức $h=5$: $P$ nằm trên đường kẻ ngang
thứ $5$. Ở độ cao $5$ (giữa chiều cao), tam giác rộng bằng một nửa đáy, từ vị trí $2{,}5$ đến $7{,}5$. Các điểm lưới
nằm \emph{bên trong} trên đoạn đó là $(3,5),(4,5),(5,5),(6,5),(7,5)$: có $\mathbf{5}$ điểm.

\step{3}{Tam giác $FPG$: dùng đường thẳng song song}
\begin{recallbox}Các tam giác có chung đáy và có đỉnh thứ ba nằm trên một đường thẳng song song với đáy thì có
cùng chiều cao, nên có diện tích bằng nhau.\end{recallbox}
Gọi $M=(5,0)$ là trung điểm của $GH$. Tam giác $FMG$ có đáy $GM=5$, chiều cao $10$, nên $S_{FMG}=\frac12\times5\times10=25$.
Vậy mọi điểm $P$ nằm trên đường thẳng đi qua $M$ và song song với $FG$ đều cho $S_{FPG}=25$.
\begin{proofbox}Nếu $P$ nằm xa $FG$ hơn đường thẳng đó thì chiều cao của tam giác $FPG$ (ứng với đáy $FG$) lớn hơn, nên
diện tích lớn hơn $25$; nếu nằm gần $FG$ hơn thì diện tích nhỏ hơn $25$. Vậy $S_{FPG}=25$ \emph{chỉ khi} $P$ nằm trên
đường thẳng ấy.\end{proofbox}
Cạnh $FG$ đi từ $(0,0)$ tới $(5,10)$: cứ sang phải $1$ ô thì lên $2$ ô. Đường song song qua $M=(5,0)$ cũng vậy:
đi qua $(5,0)$, $(6,2)$, $(7,4)$ rồi gặp cạnh $FH$ tại $(7{,}5;\,5)$ (trung điểm của $FH$). Các điểm lưới nằm
\emph{bên trong} tam giác là $(6,2)$ và $(7,4)$: có $\mathbf{2}$ điểm. ($M$ nằm trên cạnh $GH$ nên bị loại.)

\step{4}{Tam giác $HPF$: đối xứng}
Hình đối xứng qua đường thẳng đứng đi qua $F$, nên tương tự có $\mathbf{2}$ điểm: $(4,2)$ và $(3,4)$.
\begin{figC}\begin{scope}[scale=0.42]
  \draw[Line,thin] (0,0) grid (10,10);
  \draw[Ink,line width=1.2pt] (0,0)--(10,0)--(5,10)--cycle;
  \draw[tmBlue,line width=1.2pt] (2.5,5)--(7.5,5);
  \draw[tmOrange,line width=1.2pt] (5,0)--(7.5,5);\draw[tmGreen!80!black,line width=1.2pt] (5,0)--(2.5,5);
  \foreach \p in {(3,5),(4,5),(5,5),(6,5),(7,5)}{\fill[tmBlue] \p circle (5pt);}
  \foreach \p in {(6,2),(7,4)}{\fill[tmOrange] \p circle (5pt);}
  \foreach \p in {(4,2),(3,4)}{\fill[tmGreen!80!black] \p circle (5pt);}
  \node[above,font=\scriptsize] at (5,10) {$F$};\node[below left,font=\scriptsize] at (0,0) {$G$};
  \node[below right,font=\scriptsize] at (10,0) {$H$};\node[below,font=\scriptsize] at (5,0) {$M$};
  \node[right,font=\scriptsize,text=tmBlue,align=left] at (10.4,6) {xanh dương: $S_{GPH}=25$};
  \node[right,font=\scriptsize,text=tmOrange!80!black] at (10.4,4.6) {cam: $S_{FPG}=25$};
  \node[right,font=\scriptsize,text=tmGreen!70!black] at (10.4,3.2) {xanh lá: $S_{HPF}=25$};
\end{scope}\end{figC}

\step{5}{Cộng lại}
Số tam giác có diện tích $25$: $5+2+2=9$. (Một điểm $P$ không thể cho hai tam giác cùng có diện tích $25$, vì khi đó tam giác
thứ ba có diện tích $50-25-25=0$, mà $P$ nằm hẳn bên trong. Nên không có tam giác nào bị đếm thiếu hay trùng.)
\verify{Với $P=(6,2)$: tam giác $FPG$ nằm trong hình chữ nhật từ $(0,0)$ đến $(6,10)$, diện tích $60$. Cắt bỏ ba tam giác
vuông: $\frac12\times6\times2=6$, $\frac12\times1\times8=4$ và $\frac12\times5\times10=25$. Còn lại $60-6-4-25=25$. Đúng!}
Vậy chọn \textbf{(E)}.
\ans{(E) $9$ tam giác}
}

\begin{probbox}{27}{\relax}
Người ta xếp $100$ điểm thành một lưới vuông $10\times10$ như hình bên. Điểm $P$ đã được đánh dấu. Chọn ngẫu nhiên một
điểm $Q$ trong $99$ điểm còn lại. Xác suất để đoạn thẳng $PQ$ nằm ngang hoặc thẳng đứng là
\begin{figC}\begin{scope}[scale=0.38]
  \foreach \x in {0,...,9}{\foreach \y in {0,...,9}{\fill[Navy] (\x,\y) circle (3.2pt);}}
  \fill[white] (6,3) circle (6pt);\draw[tmOrange,line width=1.2pt] (6,3) circle (6pt);
  \node[above right,font=\scriptsize,text=tmOrange!80!black] at (6,3) {$P$};
\end{scope}\end{figC}
\choices{$\dfrac2{11}$}{$\dfrac15$}{$\dfrac1{10}$}{$\dfrac4{25}$}{$\dfrac5{33}$}
\end{probbox}
\probsol{
\textbf{Phân tích đề.} $Q$ được chọn ngẫu nhiên trong $99$ điểm (không tính $P$), các điểm đồng khả năng.
Đoạn $PQ$ nằm ngang khi $Q$ cùng hàng với $P$; thẳng đứng khi $Q$ cùng cột với $P$.

\step{1}{Số kết quả có thể}
Có $100-1=99$ cách chọn $Q$.

\step{2}{Số kết quả thuận lợi}
Hàng chứa $P$ có $10$ điểm, trừ $P$ còn $9$ điểm. Cột chứa $P$ có $10$ điểm, trừ $P$ còn $9$ điểm.
Một điểm khác $P$ không thể vừa cùng hàng vừa cùng cột với $P$ (điểm duy nhất như thế là chính $P$), nên
không có điểm nào bị đếm hai lần: tổng cộng $9+9=18$ điểm.
\begin{figC}\begin{scope}[scale=0.3]
  \foreach \x in {0,...,9}{\foreach \y in {0,...,9}{\fill[Slate!50] (\x,\y) circle (3pt);}}
  \foreach \x in {0,...,9}{\fill[tmBlue] (\x,3) circle (5pt);}
  \foreach \y in {0,...,9}{\fill[tmGreen!80!black] (6,\y) circle (5pt);}
  \fill[white] (6,3) circle (6pt);\draw[tmOrange,line width=1.2pt] (6,3) circle (6pt);
  \node[right,font=\scriptsize,text=tmBlue] at (10,3) {$9$ điểm cùng hàng};
  \node[above,font=\scriptsize,text=tmGreen!70!black] at (6,9.3) {$9$ điểm cùng cột};
\end{scope}\end{figC}

\step{3}{Tính xác suất}
\[P=\frac{18}{99}=\frac{18:9}{99:9}=\frac{2}{11}.\]
\begin{warnbox}Mẫu số là $99$ chứ không phải $100$, vì $Q$ phải khác $P$. Dùng $100$ sẽ ra $\frac{18}{100}$, không có trong
các phương án. Cũng đừng đếm $10+10=20$ điểm (tính cả $P$ hai lần).\end{warnbox}
\begin{tipbox}Vị trí của $P$ trên lưới không quan trọng: điểm nào cũng có đúng $9$ điểm cùng hàng và $9$ điểm cùng cột.\end{tipbox}
\verify{$\frac2{11}\times99=18$, đúng bằng số điểm thuận lợi.}
Vậy chọn \textbf{(A)}.
\ans{(A) $\dfrac2{11}$}
}

\begin{probbox}{28}{\stars{1}}
Tổng tất cả các chữ số của các số nguyên từ $98$ đến $101$ là
\[9+8+9+9+1+0+0+1+0+1=38.\]
Hỏi tổng tất cả các chữ số của các số nguyên từ $1$ đến $2008$ là bao nhiêu?
\choices{$30\,054$}{$27\,018$}{$28\,036$}{$30\,036$}{$28\,054$}
\end{probbox}
\probsol{
\textbf{Phân tích đề.} Viết ra $2008$ số rồi cộng là quá dài. Ta chia thành ba khối: từ $1$ đến $999$; từ $1000$ đến
$1999$; từ $2000$ đến $2008$, và tính tổng chữ số của từng khối bằng cách đếm \emph{mỗi chữ số xuất hiện bao nhiêu lần}.

\step{1}{Khối $1$ đến $999$: viết thêm số $0$ vào trước cho đủ ba chữ số}
Viết các số từ $0$ đến $999$ dưới dạng ba chữ số: $000, 001, 002, \dots, 998, 999$. Thêm chữ số $0$ vào trước và thêm số
$000$ \emph{không làm thay đổi} tổng các chữ số. Có tất cả $1000$ số, tức $3000$ chữ số.
\begin{proofbox}Ở hàng đơn vị, các chữ số lặp lại theo chu kì $0,1,2,\dots,9$; trong $1000$ số có $100$ chu kì nên mỗi chữ số
$0$--$9$ xuất hiện đúng $100$ lần. Ở hàng chục, mỗi chữ số đứng liền $10$ lần ($00$--$09$ có chữ số hàng chục $0$, \dots)
và lặp lại $10$ lần, cũng được $100$ lần. Ở hàng trăm: $000$--$099$ có chữ số hàng trăm $0$ ($100$ lần), \dots,
$900$--$999$ có chữ số hàng trăm $9$ ($100$ lần). Vậy ở mỗi hàng, mỗi chữ số xuất hiện $100$ lần.\end{proofbox}
Do đó mỗi chữ số $0,1,\dots,9$ xuất hiện $3\times100=300$ lần. Tổng các chữ số:
\[300\times(0+1+2+\dots+9)=300\times45=13\,500.\]
\begin{recallbox}$0+1+2+\dots+9=\dfrac{9\times10}{2}=45$ (ghép cặp $0+9=1+8=\dots=4+5=9$, có $5$ cặp).\end{recallbox}

\step{2}{Khối $1000$ đến $1999$}
Mỗi số có dạng $1\,abc$ với $abc$ chạy từ $000$ đến $999$. Phần ba chữ số cuối cho tổng $13\,500$ (như bước 1). Chữ số
hàng nghìn là $1$, xuất hiện $1000$ lần, góp thêm $1000$. Tổng: $13\,500+1000=14\,500$.

\step{3}{Khối $2000$ đến $2008$}
Có $9$ số. Chữ số hàng nghìn $2$ góp $9\times2=18$; hai chữ số giữa đều là $0$; chữ số hàng đơn vị là $0,1,\dots,8$, góp
$0+1+\dots+8=36$. Tổng: $18+36=54$.

\step{4}{Cộng ba khối}
\[13\,500+14\,500+54=28\,054.\]
\begin{warnbox}Đừng quên khối lẻ cuối cùng $2000$--$2008$ (nếu quên sẽ được $28\,000$), và đừng quên chữ số hàng nghìn $1$
của khối thứ hai (nếu quên sẽ được $27\,054$).\end{warnbox}
\verify{Thử cách làm trên một trường hợp nhỏ, từ $1$ đến $99$: mỗi chữ số xuất hiện $2\times10=20$ lần, tổng
$20\times45=900$. Đếm trực tiếp: hàng đơn vị góp $10\times45=450$, hàng chục góp $10\times(1+2+\dots+9)=450$, cộng lại $900$. Khớp.}
Vậy chọn \textbf{(E)}.
\ans{(E) $28\,054$}
}

\begin{probbox}{29}{\relax}
Gieo hai con xúc xắc cân đối, một con màu đỏ và một con màu xanh (mỗi con có sáu mặt ghi số từ $1$ đến $6$). Xác suất để
số chấm trên con xúc xắc đỏ lớn hơn số chấm trên con xúc xắc xanh là
\choices{$\dfrac{18}{36}$}{$\dfrac{25}{36}$}{$\dfrac{15}{36}$}{$\dfrac{12}{36}$}{$\dfrac{17}{36}$}
\end{probbox}
\probsol{
\textbf{Phân tích đề.} Một kết quả là một cặp (số trên con đỏ, số trên con xanh). Hai con xúc xắc khác màu nên
$(5,2)$ và $(2,5)$ là hai kết quả khác nhau.

\step{1}{Số kết quả có thể}
Theo quy tắc nhân: $6\times6=36$ kết quả, đồng khả năng.

\step{2}{Lập bảng và đánh dấu các ô ``đỏ $>$ xanh''}
\begin{figC}\begin{scope}[scale=0.5]
  \foreach \r in {1,...,6}{\foreach \b in {1,...,6}{
    \pgfmathsetmacro\yy{6-\r}
    \ifnum\r>\b \fill[AnsRed!25] (\b-1,\yy) rectangle ++(1,1);
    \else\ifnum\r=\b \fill[Slate!20] (\b-1,\yy) rectangle ++(1,1);\fi\fi}}
  \draw[Ink] (0,0) grid (6,6);
  \foreach \i in {1,...,6}{\node[font=\scriptsize,text=tmBlue] at (\i-0.5,6.4) {\i};
     \node[font=\scriptsize,text=AnsRed] at (-0.4,6.5-\i) {\i};}
  \node[font=\scriptsize,text=tmBlue] at (3,7.1) {xúc xắc xanh};
  \node[font=\scriptsize,text=AnsRed,rotate=90] at (-1.2,3) {xúc xắc đỏ};
\end{scope}\end{figC}
Ô tô đỏ: số đỏ lớn hơn số xanh; ô tô xám: hai số bằng nhau. Đếm theo từng hàng (số đỏ):
đỏ $=1$: $0$ ô; đỏ $=2$: $1$ ô; đỏ $=3$: $2$ ô; đỏ $=4$: $3$ ô; đỏ $=5$: $4$ ô; đỏ $=6$: $5$ ô.
Tổng: $0+1+2+3+4+5=15$ kết quả thuận lợi.

\step{3}{Tính xác suất}
$P(\text{đỏ}>\text{xanh})=\dfrac{15}{36}$.

\step{4}{Cách nhanh: dùng tính đối xứng}
Có $6$ kết quả hai số bằng nhau (đường chéo). Còn lại $36-6=30$ kết quả, trong đó ``đỏ lớn hơn'' và ``xanh lớn hơn'' có số
lượng bằng nhau (đổi vai hai con xúc xắc), nên mỗi loại có $30:2=15$ kết quả.
\begin{warnbox}Đáp án $\frac{18}{36}=\frac12$ là bẫy: nó quên mất các trường hợp hai số bằng nhau. ``Đỏ lớn hơn'' và
``xanh lớn hơn'' không phải là hai khả năng duy nhất.\end{warnbox}
\verify{$P(\text{đỏ}>\text{xanh})+P(\text{bằng nhau})+P(\text{xanh}>\text{đỏ})=\frac{15}{36}+\frac6{36}+\frac{15}{36}=1$.}
Vậy chọn \textbf{(C)}.
\ans{(C) $\dfrac{15}{36}$}
}

\begin{probbox}{30}{\relax}
Gieo hai con xúc xắc cân đối, một con màu đỏ và một con màu xanh. Xác suất để \emph{tích} hai số chấm xuất hiện là một số
chẵn bằng
\choices{$\dfrac12$}{$\dfrac34$}{$\dfrac14$}{$\dfrac23$}{$\dfrac{11}{36}$}
\end{probbox}
\probsol{
\textbf{Phân tích đề.} Như bài 29, có $6\times6=36$ kết quả đồng khả năng. Đếm trực tiếp các tích chẵn khá dài, nên ta
đếm \emph{phần bù}: các kết quả có tích lẻ.

\step{1}{Khi nào tích hai số là số lẻ?}
\begin{recallbox}lẻ $\times$ lẻ $=$ lẻ;\quad chẵn $\times$ (số bất kì) $=$ chẵn. Vậy tích hai số là số lẻ khi và chỉ khi
\emph{cả hai} số đều lẻ.\end{recallbox}

\step{2}{Đếm các kết quả có tích lẻ}
Số lẻ trên mỗi con xúc xắc: $1,3,5$ ($3$ số). Cả hai con đều lẻ: $3\times3=9$ kết quả.

\step{3}{Đếm các kết quả có tích chẵn bằng phần bù}
\begin{recallbox}\emph{Phần bù}: số kết quả thuận lợi $=$ số kết quả có thể $-$ số kết quả không thuận lợi. Tương tự,
$P(\text{biến cố})=1-P(\text{biến cố không xảy ra})$.\end{recallbox}
Số kết quả có tích chẵn: $36-9=27$.

\step{4}{Tính xác suất}
\[P(\text{tích chẵn})=\frac{27}{36}=\frac{27:9}{36:9}=\frac34.\]
\begin{warnbox}Nhiều bạn nghĩ ``chẵn hay lẻ là nửa nọ nửa kia'' nên chọn $\frac12$. Nhưng chỉ cần \emph{một} số chẵn là tích đã
chẵn, nên tích chẵn xảy ra nhiều hơn hẳn.\end{warnbox}
\verify{Đếm trực tiếp: con đỏ chẵn ($3$ cách) với con xanh bất kì ($6$ cách): $18$ kết quả; con đỏ lẻ ($3$ cách) và con xanh
chẵn ($3$ cách): $9$ kết quả. Tổng $18+9=27$. Khớp.}
Vậy chọn \textbf{(B)}.
\ans{(B) $\dfrac34$}
}
